\documentclass[../ICIT-Paper-2024.tex]{subfiles}
\usepackage{array} % Để tạo ra các cột có định dạng cố định
\usepackage{booktabs} % Để sử dụng các đường kẻ đẹp trong bảng
\usepackage{numprint}
\usepackage{tabularx}
\begin{document}

\subsection{Dataset}

\subsubsection{Data Source}
\paragraph{}
In DeFi, data is stored and easily collected on blockchain networks. Currently, the most popular blockchain networks are those with the same technology and mechanism as Ethereum~\footnote{\href{https://ethereum.org/}{https://ethereum.org}}, such as Binance Smart Chain(BSC)~\footnote{\href{https://www.bnbchain.org/en/bnb-smart-chain}{https://www.bnbchain.org}}, Polygon~\footnote{\href{https://polygon.technology/}{https://polygon.technology}}, Tron~\footnote{\href{https://tron.network/}{https://tron.network/}}, Arbitrum~\footnote{\href{https://arbitrum.io/}{https://arbitrum.io}}, and Optimism~\footnote{\href{https://optimism.io/}{https://optimism.io}}, etc. Ethereum is a protocol for executing a virtual computer in an open and distributed manner. This virtual computer is called the Ethereum Virtual Machine (EVM). As of May 2024, the total value locked (TVL) in EVM chains accounts for over 80\% (TVL in the Ethereum chain accounts for over 60\%) of the total TVL in DeFi. We collect and aggregate multichain data from these EVM chains.

Data collection is implemented according to the model of the paper~\citeauthor{pham2022analysis}~\cite{pham2022analysis}. The collected data consists of two types: (1) raw data, which includes transactions and events, and (2) state data of entities on EVM chains, such as the balance of wallets and the price of tokens, etc. After collecting the data, we process and build a multichain knowledge graph. This knowledge graph includes two main components: (1) data of entities on the web3 space, such as wallets, tokens, projects, and dapps, and (2) relationships among entities, such as deposit, borrow, withdraw, and repay of wallets with lending protocols.

As of May 2024, we have collected 20TB of transaction and event data across 7 EVM chains: BSC, Ethereum, Avalanche~\footnote{\href{https://www.avax.network/}{https://www.avax.network}}, Polygon, Optimism, Arbitrum, and Fantom~\footnote{\href{https://fantom.foundation/}{https://fantom.foundation/}}. From this data, we have aggregated the knowledge graph data of over 600 million wallets, 25 million smart contracts, \numprint{20000} tokens, and \numprint{46000} projects. From this knowledge graph, we have selected multichain data of \numprint{251290} wallets \footnote{\href{https://github.com/sajol6436/Defi-Credit-Score/tree/main/data}{https://github.com/sajol6436/Defi-Credit-Score/tree/main/data}} to serve the parameter optimization problem for the Credit Score. These \numprint{251290} wallets were selected based on two factors: (1) total assets of the wallet over \$10 and (2) wallets that have engaged in transactions within the past year on the two main lending pools, Aave and Compound. The table~\ref{tab1} shows the parameters in the scoring formula and the corresponding feature data stored in wallets.
% \renewcommand{\arraystretch}{1.5}
\begin{table}[htbp]
    \centering
    \caption{Parameters in the scoring formula and the corresponding feature data stored in wallets.}\label{tab1} % Chú thích của bảng
    \begin{tabular}{|l|m{8cm}|} % Các cột được căn lề trái (l) và cột thứ hai có chiều rộng tối đa là 8cm
        \hline
        \textbf{Parameters} & \textbf{Description} \\ % Tiêu đề của từng cột
        \hline
        total current asset & total current asset \\ \hline
        average total asset & average total assets in the last 30 days \\ \hline
        frequency of DApp transactions & represents the activeness of a wallet on blockchain investment channels. \\ \hline
        number of interacted DApps & represents the diversity of blockchain investment made by the wallet. \\ \hline
        types of interacted DApps & the total number of DApp types that a wallet transacts within the last 30 days. \\ \hline
        reputation of interacted DApps & the number of high reputation DApps that the wallet interacted with in the last k days. \\ \hline
        transaction amount & total amount in US dollars that other wallets transfer into the wallet in the last 30 days. \\ \hline
        frequency of transaction & Frequency of transaction assesses the number of transactions that the account has performed in the last 30 days. \\ \hline
        age of accounts & denotes how long the user has participated in the crypto market, since their first transaction \\ \hline
        number of liquidations & number of liquidations \\ \hline
        total value of liquidations & total value of liquidations \\ \hline
        loan-to-balance ratio & indicates the quantitative relationship between the total loans and the balance of an account.  \\ \hline
        loan-to-investment ratio & the number of penalty points that wallet addresses will be minus due to their illegal “loan-to-investment ratio”. \\ \hline
        investment-to-total-asset ratio & represents the activeness of the account in the crypto business. \\ \hline
    \end{tabular}
    \label{tab:parameters} % Nhãn của bảng để tham chiếu sau này
\end{table}

\end{document}
